\documentclass[11pt]{article}
\usepackage[UTF8]{ctex} 
\usepackage[margin=1.1in]{geometry}
\usepackage{booktabs}
\usepackage{xcolor}
\usepackage{amsmath}

\title{\textbf{Buoyancis: 信任安全网协议 (SNP)}}
\author{创始合伙人机密简报}
\date{2026年1月}

\begin{document}

\maketitle

\section{核心愿景：跨文化信任桥梁}
Buoyancis 不仅仅是一个技术协议，它是首个通过算法确保“东西方利益对等”的信用平台。在当前全球碎片化的环境下，我们通过数学（而非仅靠法律）来重建商业信任。

\section{50/50 全球合伙人架构}
为了实现中立性，我们采取了独特的股权对等策略：

\begin{center}
\renewcommand{\arraystretch}{1.3}
\begin{tabular}{lcp{6cm}}
    \toprule
    \textbf{分部 (Division)} & \textbf{股权配额} & \textbf{战略职能} \\
    \midrule
    西方市场 (Western) & 5.625\% & 硅谷合规、全球定价权、技术标准 \\
    东方市场 (Eastern) & 5.625\% & 供应链整合、亚洲市场渗透、高并发应用 \\
    \midrule
    \textbf{总计出让} & \textbf{11.25\%} & \textbf{投后估值：\$5,000,000} \\
    \bottomrule
\end{tabular}
\end{center}

\section{合伙人 Offer (Founding Partner)}
我们为 Rick, Donnie, Johan 以及东方核心伙伴预留了专属席位（基于 5.625\% 区域池均分）：
\begin{itemize}
    \item \textbf{个人持股：} 1.875\% 原始股权
    \item \textbf{名义价值：} \$93,750
    \item \textbf{核心机制：} 遵循“安全网协议”的释放逻辑（Vesting），确保长期绑定。
\end{itemize}

\section{信任方程 ($R$)}
\[ R = \frac{\sum (M_{v} \times S_{t})}{C_{a}} \]
$R$ 代表可靠性。无论在纽约还是上海，这套公式都将作为衡量项目风险的唯一准则。

\end{document}